\section{How Part 11 was run}\label{app:provenance}

Part 11 and this paper were produced by an AI system, Claude (Anthropic), on 2 and 3 November 2026, in one orchestrating session that delegated parts of the work to separate agents. Kaarel H\"anni posed the request (\cref{app:choice}). No human has reviewed the code, the results or this text. Times are UTC, and the commits named are on the working branch \texttt{claude/nifty-mccarthy-j0lw2a}. The records are:
\begin{itemize}
\item the protocol \repo{study/part11/PROTOCOL.md}, with its post-hoc addendum;
\item the runners \repo{study/part11/eval_dfg.py} (toy suites) and \repo{study/part11/dfg_real.py} (tabular data), the shard merger \repo{study/part11/merge.py}, and the comparison script \repo{study/part11/compare_dfg.py} with its outputs \repo{study/part11/RESULTS.md} and \repo{study/part11/RESULTS.json};
\item the five analyses in \repo{study/part11/interp/}, each with its scripts, tables and result files, and the verification records of four of them (\cref{tab:verif});
\item this paper's generators and check scripts in \repo{paper_dfg/}.
\end{itemize}

\paragraph{Organisation.} The orchestrating session chose the object, wrote the protocol and the comparison script, and ran the queue of pre-registered runs and the core. It delegated three kinds of work to separate agents, each with its own instructions: five analysts (the structure atlas, the order of acquisition, the theory notes, false admissions and compute, and the ablations); one independent verifier for each analysis; and writing agents for the sections of this paper. Each writing agent drafted from the records it names. The abstract, the introduction, the discussion and this appendix were written last, from the drafted sections and the records.

% ------------------------------------------------------------------------------------------------------------------
\subsection{The request and the choice of object}\label{app:choice}

After the Part 10 paper \cite{portfoliolearner2026}, the request was: ``can you please suggest a discrete method that is more like `one thing', less an ensemble of different things. eg maybe one of the components of your current proposal would do? the aim behind this last question is to get a specific interesting object worth studying. then write an analogous paper but specifically about that''. The orchestrating session chose DFG v1.3 run alone, for the three reasons of \cref{sec:intro}, and set aside v3's image boosters, the portfolio itself and a new DFG for real-valued data. It wrote the choice, its reasons and its domain limit into the protocol before any Part 11 run. The choice was made after every Part 10 result had been read, v3's on \code{dlheld7} included, and the protocol says so (\cref{sec:part11}).

% ------------------------------------------------------------------------------------------------------------------
\subsection{The protocol and the runs}\label{app:runs}

\paragraph{The protocol} was written on 2 November and committed at 22:15:39 together with the runners and the merge script (\texttt{aca1f3c8}), before the first run. It fixed the object, the seven questions with their runs, suites and seeds, the rules of Parts 8--10, the rule for crashes, and what is held out (\cref{sec:part11}). Its pre-registered part has not changed since. The addendum that defines the core run was committed at 02:25:54 on 3 November (\texttt{12eae4e0}; its heading says 02:30), after every pre-registered run had finished and its results had been read, and two seconds before the core's first job started.

\paragraph{The queue.} The pre-registered runs were 282 jobs in one file:
\begin{itemize}
\item 48 shards of DFG alone on the 11 toy suites, 3 to 7 per suite;
\item 11 \code{lowdeg} jobs, one per suite;
\item 18 jobs at $d=15$, 60 and 120, one per target and $d$;
\item one job for the two tabular datasets;
\item 204 ablation shards, 17 for each of the 12 ablations over the five ablation suites.
\end{itemize}
Four worker processes started at 22:15:41. Each worker took the next unclaimed job in file order (claiming it by creating a directory, which is atomic), ran it at nice 10 with single-threaded BLAS, and marked it done or failed. A merger joined the shards of each run and suite with \texttt{merge.py} once all of them were done, and the orchestrator committed the merged files. No job failed. The first results, DFG alone on the held-out \code{dlheld7}, were committed at 22:23 (\texttt{5bd788c4}). DFG alone had finished on all 11 suites by 23:31, and \code{lowdeg}, the runs at other $d$ and the tabular runs by 23:41. The ablations' results were committed from 23:42 on, and at 02:24 all 282 jobs, 87 result files, were complete (\texttt{28d64a85}).

\paragraph{The container restart} at about 01:09 on 3 November stopped every process. Four ablation shards were running: \code{noblocks} on \code{dlheld6} (shards 2 and 3) and on \code{dlheld7} (shards 1 and 2). A shard writes its result file only when it finishes, so these four had written nothing. Their claims were removed, three workers were restarted at 01:10:54, and the four shards ran again from the start, with the same code, seeds and random streams. The 158 jobs already finished kept their files, and the merger skips any run and suite merged before a restart. So no finished result was lost or computed twice, and no result file mixes two runs; the worker logs show 286 starts for the 282 jobs. After the restart the queue ran on three workers instead of four. The calibration analysis's record writer was also stopped, at 8 of the 220 noise fits at $d=60$ and 32 of the 105 profile records. It was resumed from its records, which skip every key already written, so no cell was fitted twice. The commit \texttt{ba81503c} (01:10:57) kept the analysts' partial edits.

\paragraph{The post-hoc runs.}
\begin{itemize}
\item \emph{The core} (\cref{sec:core}): 36 shards on four workers, from 02:25:56 to 02:38:32; its 11 result files were committed by 02:39 (\texttt{22e45319}). With them, the 98 result files of Part 11 are all complete (\repo{study/part11/RESULTS.md}, ``Run status'').
\item \emph{The paired null} (\repo{study/part11/interp/ablations/perm_null.py}) on all 11 suites, seeds 0 and 1, after the core had finished and until 03:36; \emph{the leave-one-in runs} (\repo{study/part11/interp/ablations/leavein.py}), until 03:24; and 36 refits of single cells for the ablation analysis (\repo{study/part11/interp/ablations/refit.py}). Every refit reproduces its run's result file in test error and FLOP count. The leave-one-in targets were chosen after reading the core's results, and five of its configurations after reading the first leave-one-in results (they are marked in \texttt{leavein.py}).
\item \emph{This paper's checks}: \repo{paper_dfg/check_parity_variant.py} (43 refits of the parity cells at small $n$) and \repo{paper_dfg/check_h5s1.py} (6 refits of h5s1). Every refitted test error equals the result file's.
\end{itemize}
All of these ran single-threaded, as separate processes outside the queue.

% ------------------------------------------------------------------------------------------------------------------
\subsection{The analyses and their verification}\label{app:interp-workflow}

\paragraph{The analyses} ran from about 23:00 on 2 November, in parallel with the queue, on the result files as they landed (\cref{sec:interp}). Each fit they analyse is an instrumented rerun of an evaluation fit, and each analysis checks its reruns against the result files in test error and FLOP count (\cref{app:interp-detail}). The theory notes' proofs are about the frozen code, line by line. The protocol fixed Q7's questions, not the analyses' methods: each analysis chose its targets, definitions and groupings itself, and the atlas revised some definitions after its first fits had been looked at (\cref{sec:interp-how}). The paired null, the leave-one-in runs, the refits and the target families of the ablation analysis are post hoc (\cref{sec:ablations}).

\begin{table}[ht]
\centering\footnotesize\setlength{\tabcolsep}{4pt}
\begin{tabularx}{\linewidth}{@{}ll>{\raggedright\arraybackslash}Xcccc@{}}
\toprule
report & verification & what it covers & \hd{claims\\checked} & confirmed & corrected & \hd{omissions\\listed}\\
\midrule
\texttt{ATLAS.md} & \texttt{VERIFY\_ATLAS} & what the models contain: 37 targets, 407 fits & 61 & 39 & 22 & 9\\
\texttt{ORDER.md} & \texttt{VERIFY\_ORDER} & the order of acquisition, against the Part 5 networks & 60 & 41 & 19 & 9\\
\texttt{THEORY.md} & \texttt{VERIFY\_THEORY} & the exact statements P1--P5 and their checks & 54 & 47 & 7 & 4\\
\texttt{CALIB.md} & \texttt{VERIFY\_CALIB} & false admissions; where the compute goes & 28 & 16 & 12 & 8\\
\texttt{ABLATIONS.md} & -- & the ablations and the core DFG & \multicolumn{4}{c}{not verified when this table was made}\\
\bottomrule
\end{tabularx}

\caption{The independent verifications of the Part 11 analyses, counted from the verifiers' own records \status{Numerical}. ``Corrected'': the verifier changed a number, narrowed a statement or rejected a reading; ``omissions'': further problems it listed that were not tied to one claim. The reports are in \repo{study/part11/interp/}, each next to its verification record (\texttt{.md} and \texttt{.json}). Generated by \repo{paper_dfg/make_tables.py} (\texttt{tab\_verif}) from the \texttt{VERIFY\_*.json} files.}
\label{tab:verif}
\end{table}

\paragraph{The verifications.} Each verifier was a separate agent with its own instructions. It read the report, the code and the stored files, recomputed the numbers with its own scripts, refitted cells through its own code path, and listed every claim as confirmed or corrected, with the evidence. The analyst then checked each correction against the stored files, applied it or gave its reason for declining, and listed every change in the report's Verification section. \Cref{tab:verif} counts the verdicts. The paper uses the corrected numbers and wording throughout.
\begin{itemize}
\item \emph{The structure atlas} (\texttt{ATLAS.md}; verified at about 00:36). The verifier refitted 11 cells and 6 counterfactual variants and recomputed every count. All 22 corrections were applied and none was declined. Among them: the 37 targets come from 9 of the 11 suites; block representations are checked only as functions of the block statistics; thrint is first fitted without error at $n=512$ by a product block, not at 1024 by weighted blocks; and DFG alone's model of the residue classes is constant at every $n$, whereas v3's member was not at small $n$.
\item \emph{The order of acquisition} (\texttt{ORDER.md}; about 00:34). The verifier made 12 refits through a separate code path and recomputed all 23,727 round-0 $z$-scores. All its corrections and omissions were applied, two of them only in part, with the stored evidence for the rest, and the analyst's own re-check found eight further errors (for example 27 fits with $n\ge64$, not 30). Among them: DFG's round order is negative in 3 of the 69 multi-round fits; the $n$-axis comparison is like for like in data and rule, but not in the precision of the networks' estimate for small coefficients; and DFG's bundles are cut by detectability, not by coefficient size alone.
\item \emph{The theory notes} (\texttt{THEORY.md}; about 02:59). The verifier re-derived every proof against the cited lines of the code, refitted 615 cells, and gave the verdict ``sound''. All 7 corrections and 4 omissions were applied. No proof or measured value changed. The corrections narrowed wording that overstated what is proved (passing a test against being admitted to the model; $n_{\rm thr}$'s exact decomposition) and fixed one count, and the omissions added P1a$'$, the bound for other input laws (\cref{sec:guarantees}).
\item \emph{False admissions and compute} (\texttt{CALIB.md}; about 02:43). The verifier rebuilt the instrumented copy byte for byte, showed that its syntax tree minus the hooks equals the frozen file's, and made 17 independent refits. It corrected 12 claims, among them the identity claim (only the checked fits are checked; the others equal the frozen learner by construction), the linear block's null rate (close to the $\chi^2_d$ tail, not at it), sum-block admissions in the fold fits, a count (229, not 238) and one step of a proof. Its important omission is the spectral block's null pass rate on biased labels, which grows with $d$ to 99\% at $d=120$. \Cref{sec:interp-calib} reports the corrected rates.
\item \emph{The ablations} (\texttt{ABLATIONS.md}) were being checked by an independent verifier while this paper was completed. \Cref{sec:ablations} reports the analysis as written.
\end{itemize}

% ------------------------------------------------------------------------------------------------------------------
\subsection{Disclosures}\label{app:disclosures}

\begin{itemize}
\item \emph{The timing check.} Before the protocol was written, DFG alone was fitted on two \code{dlheld7} targets, f7\_03 and f7\_08 (seed 0; $n=64$, 512, 4096 and 16384), to time it, and the 8 test errors were printed. The protocol records them, and the queue's runs reproduce all 8 to the three decimals recorded (\repo{study/part11/RESULTS.md}, ``Consistency checks''). The code had been frozen in Part 10 and nothing is tuned in Part 11.
\item \emph{Smoke tests} of the ablations and of \code{lowdeg} used three \code{designed} targets: maj9, stair5 and par4.
\item \emph{A replaced setting.} The first choice for \code{lenient}, $\delta=10$, crashed: the linear block's threshold takes the square root of a multiple of $\ln(1/\delta)$, which is negative for $\delta>1$. It was replaced by $\delta=0.9$ before any run, and the protocol records both.
\item \emph{The object was chosen after v3's results} on every suite had been read, \code{dlheld7}'s included (\cref{app:choice}).
\item \emph{Post-hoc choices.} The core was defined after every pre-registered result had been read, and the leave-one-in targets after the core's (\cref{sec:core}). The target families and classes were assigned by hand after the results were in.
\item \emph{Load.} The queue ran four workers at nice 10 on a four-CPU machine, three after the restart, next to the analyses, so the CPU times are rough (\cref{sec:res-compute}).
\item \emph{Written and checked by AI agents.} The analyses, their verifications and this paper were produced by AI agents (Claude), each with its own instructions. No human has reviewed them.
\end{itemize}

% ------------------------------------------------------------------------------------------------------------------
\subsection{Reproducing the numbers}\label{app:repro}

\paragraph{Which files stand for which learner.}
\begin{itemize}\raggedright
\item \emph{DFG alone, its ablations, the core and \code{lowdeg}} on the toy suites: \texttt{toy\_\_dfg13-perm\_\_<suite>.json}, \texttt{toy\_\_dfg13-<tag>-perm\_\_<suite>.json}, \texttt{toy\_\_dfg13-core-perm\_\_<suite>.json} and \texttt{toy\_\_lowdeg-perm\_\_<suite>.json}, all in \repo{study/part11/results/}. The paired null's files are in \repo{study/part11/interp/ablations/results/}.
\item \emph{DFG alone on tabular data}: \texttt{dfg13\_\_kr\_vs\_kp.json} and \texttt{dfg13\_\_mushroom.json} in \repo{study/standard_tasks/results/}.
\item \emph{v3 on toy inputs}: \code{p10_sum}'s files in \repo{study/hybrid/results/}, which is v3's code there (\cref{sec:refs}). The networks and the other references are read from their own result files, as the docstring of \repo{study/part11/compare_dfg.py} lists them.
\end{itemize}

\paragraph{Commands.} {\raggedright The commands are run from the folder \texttt{gd-circuit-growth/}, with single-threaded BLAS (the environment variables \texttt{OMP\_NUM\_THREADS} and \texttt{OPENBLAS\_NUM\_THREADS} set to 1):\par}
\begin{itemize}\raggedright
\item the runs: \texttt{cd study/part11}, then \texttt{python3 eval\_dfg.py dfg13 <suite>} (an ablation adds \texttt{-{}-kw=<json> -{}-name=<tag>}, with the arguments of \cref{sec:part11} and, for the core, those of the protocol's addendum; shards add \texttt{-{}-fns=} and \texttt{-{}-part=}, and \texttt{python3 merge.py <run> <suite>} joins them), \texttt{python3 eval\_dfg.py lowdeg <suite>}, \texttt{python3 eval\_dfg.py dfg13 designed\_d<d>}, and \texttt{python3 dfg\_real.py kr\_vs\_kp,mushroom};
\item the comparisons: \texttt{python3 compare\_dfg.py} (writes \texttt{RESULTS.md} and \texttt{RESULTS.json});
\item the analyses: the scripts that each report names in \texttt{interp/} (for example \texttt{atlas\_run.py} and \texttt{atlas\_make.py}; \texttt{order/order\_fit.py} and \texttt{order\_analysis.py}; \texttt{theory/theory\_fits.py} and \texttt{check\_theory.py}; \texttt{calib/calib.py} and \texttt{calib\_summary.py}; \texttt{ablations/perm\_null.py}, \texttt{leavein.py}, \texttt{refit.py} and \texttt{ablations.py});
\item this paper: \texttt{cd ../../paper\_dfg}, then \texttt{python3 make\_tables.py}, \texttt{python3 make\_figures.py}, \texttt{python3 make\_interp.py}, \texttt{python3 check\_parity\_variant.py} and \texttt{python3 check\_h5s1.py}, and \texttt{latexmk -pdf main.tex}; \texttt{dfg-paper.pdf} is a copy of \texttt{main.pdf}.
\end{itemize}
Every table and figure is generated from these files; derived numbers that the text quotes and no table shows are in \repo{paper_dfg/tables/numbers.json}, \repo{paper_dfg/tables/interp_numbers.json} and the figures' \texttt{.json} files.

\paragraph{What is not in the repository, and what stays on the working branch.} The job queue, the worker and merger scripts and their logs stayed in the session scratch directory and are in neither branch; they only wrap the committed runners, whose commands are above. The networks' raw outputs also live in a scratch directory, as in Parts 1--10. On the working branch, the LaTeX build output in \texttt{paper\_dfg/} (\texttt{main.pdf} and the auxiliary files) and the run logs of the post-hoc runs (\texttt{study/part11/interp/ablations/logs/}) stay there; \texttt{dfg-paper.pdf} is the published copy of the paper.
